\chapter{Confronto sperimentale PIR vs mmWave}
\label{cap:confronto}

\begin{quote}\small
Questo è il capitolo centrale della tesi e risponde all'obiettivo O3. Presenta il
protocollo di misura, i risultati acquisiti e la loro discussione critica. Le
sezioni contrassegnate come \emph{in corso} indicano acquisizioni pianificate ma
non ancora completate alla data di questa stesura --- penetrazione degli ostacoli,
prova del respiro a ritmo imposto, confronto quantitativo con il LD2420 e misura
angolare: la loro collocazione nel capitolo è già definita, così come le metriche e
i comandi con cui verranno prodotte.
\end{quote}

\section{Protocollo sperimentale}
\label{sec:protocollo}

\subsection{Principi generali}

Il protocollo è stato definito \emph{prima} di iniziare le acquisizioni ed è
rimasto invariato, per una ragione precisa: la campagna completa richiede circa
dieci ore di acquisizioni, e una modifica del protocollo a metà strada le renderebbe
non confrontabili.

I principi adottati sono cinque.

\begin{enumerate}[leftmargin=1.8em,itemsep=4pt]
  \item \textbf{Cinque ripetizioni per scenario} (T01--T05). È ciò che permette di
        riportare media e deviazione standard, e quindi di rendere il confronto
        ``numerico'' e non anecdottico. Ogni valore riportato in questo capitolo
        nella forma $\mu \pm \sigma$ è calcolato su cinque trial indipendenti.
  \item \textbf{Acquisizione sincrona dei due sensori}. Radar e PIR sono letti nello
        stesso frame dal medesimo firmware, sulla stessa base dei tempi: il
        confronto non richiede allineamento a posteriori e non introduce errori di
        sincronizzazione.
  \item \textbf{Ground truth dichiarata}. Ogni acquisizione porta nei metadati la
        presenza reale e lo stato reale del soggetto, dichiarati dall'operatore da
        riga di comando. Per la ground truth statica per file questo è sufficiente;
        per le latenze si adotta la convenzione dell'evento a istante noto
        (\cref{sec:latenza}).
  \item \textbf{Setup fisso}. LD2410B e PIR affiancati sul bordo di un tavolo a
        circa 1~m di altezza, puntati verso l'area di prova, con la zona di prova
        interna al campo di \emph{entrambi} i sensori (il PIR copre un cono più
        ampio, $<$110\degree, del radar). Nessun oggetto in movimento nella stanza:
        nessun ventilatore, tende ferme, porta chiusa.
  \item \textbf{Alimentazione da alimentatore di rete} da almeno 1~A e non dalla
        porta USB del portatile. L'ESP32 con WiFi attivo e il LD2410B presentano
        picchi superiori a 400~mA, e una porta USB debole provoca brownout, cioè
        riavvii casuali che si presentano come bug del firmware.
\end{enumerate}

\subsection{Metriche}

Le metriche calcolate da \texttt{analisi/analizza\_test.py} sono le seguenti.

\begin{description}[leftmargin=1.4em,style=nextline,itemsep=3pt]
  \item[Falsi negativi (\%)] percentuale dei campioni con presenza reale in cui il
    sensore riporta assenza. È la metrica chiave dello scenario ``persona ferma''.
  \item[Falsi positivi (eventi/h)] numero di transizioni $0 \to 1$ per ora di
    osservazione a stanza vuota. Si conta per \emph{eventi} e non per campioni,
    perché un singolo falso positivo che dura dieci secondi non è dieci volte più
    grave di uno che dura un secondo.
  \item[Latenza di rilevamento] tempo fra l'evento reale e la prima transizione
    $0 \to 1$ del sensore.
  \item[Latenza di rilascio] tempo fra l'uscita del soggetto e la transizione
    stabile $1 \to 0$, con conteggio separato delle eventuali riaccensioni nella
    coda.
  \item[Accuratezza della distanza] scostamento fra distanza riportata dal radar e
    distanza reale misurata, con separazione fra dispersione \emph{fra} trial e
    dispersione \emph{entro} il trial.
\end{description}

\subsection{Registro delle sessioni}
\label{sec:registro}

Ogni sessione è annotata in un registro con data, ora, temperatura della stanza,
soggetto, alimentazione, test svolti e note. La temperatura è registrata perché il
PIR degrada al crescere della temperatura ambiente (\cref{sec:pir-temperatura}); le
sessioni di questa campagna si collocano fra 25 e 29\,\celsius, condizione
sfavorevole al PIR che va tenuta presente nella discussione (\cref{sec:limiti}).

Le acquisizioni completate sono riassunte nella \cref{tab:sessioni}.

\begin{table*}[htbp]
\centering
\caption{Sessioni di acquisizione completate. La colonna ``PIR'' indica la posizione
del ponticello di ritrigger dell'\pir\ (\cref{sec:jumper}).}
\label{tab:sessioni}
\small
\begin{tabularx}{\linewidth}{@{}llrcX@{}}
\toprule
\textbf{Data} & \textbf{Scenario} & \textbf{Durata} & \textbf{PIR} & \textbf{Note} \\
\midrule
18/08/2026 & validazione logger & 134~s & L & 673 campioni, 5,00~Hz esatti \\
18/08/2026 & pilota PIR vs radar & 339,6~s & L & 1699 campioni, PIR su GPIO34 \\
18/08/2026 & pilota respiro & 91~s & L & soggetto a 40~cm, nessuna ground truth \\
18/08/2026 & stanza vuota (giorno) & 1742~s & L & 8711 campioni, 29\,\celsius \\
18/08/2026 & fermo seduto T01--T05 & 5 $\times$ 302~s & L & 7554 campioni, 2,30~m, 27\,\celsius \\
20/08/2026 & distanza 1--5~m, 25 trial & 25 $\times$ 62~s & L & cammino sul posto \\
20--21/08/2026 & stanza vuota (notte) & 23\,640,6~s & L & 118\,204 campioni, 6,57~h \\
22/08/2026 & sotto il banco, 2 serie & 10 $\times$ 302~s & L & immobile e con micro-movimenti \\
23/08/2026 & verifica ponticello H/L & 3 prove brevi & H & prova della mano, \cref{sec:jumper} \\
23/08/2026 & dose-risposta a 1~m, 3 serie & 15 trial & H & fermo, micro-movimenti, cammino \\
23/08/2026 & movimento a 2~m & 5 $\times$ 62~s & H & controllo della portata \\
23/08/2026 & sotto il banco, 2 serie ripetute & 10 $\times$ 302~s & H & coppia appaiata definitiva \\
23/08/2026 & ingresso (latenza) & 10 $\times$ 62~s & H & evento annunciato a $t=30$~s \\
23/08/2026 & uscita (rilascio) & 5 $\times$ 102~s & H & evento annunciato a $t=30$~s \\
24/08/2026 & due persone, 3 geometrie & 9 $\times$ 122~s & H & più un trial di controllo \\
24--25/08/2026 & selettività spaziale & 12 $\times$ 182~s & H & gate massimo 2 (150~cm) \\
\bottomrule
\end{tabularx}
\end{table*}

\paragraph{Nota su una sessione troncata.} La sessione notturna era programmata per
sette ore ed è terminata a 6,57~h perché un aggiornamento automatico di Windows ha
riavviato il computer. I dati restano validi: l'intervallo di campionamento è di
200~ms esatti su tutti i 118\,203 intervalli, senza buchi né campioni persi ---
l'acquisizione è semplicemente troncata. L'episodio è però istruttivo e ha prodotto
una regola di protocollo: per le acquisizioni oltre l'ora, sospendere gli
aggiornamenti di Windows e verificare che sospensione e ibernazione siano disattivate.

\section{Falsi positivi: stanza vuota}
\label{sec:falsi-positivi}

\paragraph{Procedura.} Acquisizione avviata con l'operatore che esce immediatamente
dalla stanza; scarto dei primi 60~s (stabilizzazione del PIR e uscita
dell'operatore); conteggio delle transizioni $0 \to 1$ sul resto.

\paragraph{Risultato.} Su un totale di \textbf{7,05~h} di stanza vuota
(6,57~h notturni più 0,48~h diurni), \textbf{né il radar né il PIR hanno prodotto
alcun falso positivo}. Nella sessione notturna gli unici 96 campioni con presenza
sono compresi fra $t=0$ e $t=19$~s e corrispondono all'operatore che esce dalla
stanza.

Con zero eventi osservati non è corretto scrivere ``0 eventi/h'': si riporta il
limite superiore, che dipende dal tempo di osservazione. Applicando la regola del
tre ($3/T$, che approssima il limite superiore al 95\,\% di confidenza per un
processo di Poisson con zero eventi osservati):

\begin{equation}
\lambda_{95\%} \leq \frac{3}{7{,}05~\text{h}} = 0{,}43~\text{eventi/h}
\label{eq:regola-tre}
\end{equation}

Il risultato da riportare è quindi: \textbf{falsi positivi $\leq 0{,}43$ eventi/h
per entrambi i sensori, al 95\,\% di confidenza}. Il valore è coerente con l'ordine
di grandezza indicato dalla letteratura tecnica per i moduli mmWave consumer
($<$0,05 eventi/h), che però richiederebbe una sessione di decine di ore per essere
confermato con questa metodologia. Questo è il modo corretto di riportare
un'assenza di eventi, e la principale limitazione del risultato è il tempo di
osservazione, non il sensore.

\paragraph{Nota sulla temperatura.} Le sessioni sono state svolte a 29\,\celsius. È
una condizione \emph{sfavorevole} al PIR per il rilevamento (minore contrasto
termico) ma \emph{favorevole} sui falsi positivi: un ambiente caldo e uniforme
produce meno gradienti spuri. Il risultato ``zero falsi positivi PIR'' va quindi
letto come non pessimistico.

\section{Il risultato centrale: la persona ferma}
\label{sec:persona-ferma}

\paragraph{Procedura.} Soggetto seduto immobile a 2,30~m dalla coppia di sensori,
respirazione spontanea, nessun movimento volontario. Cinque trial da 302~s utili
ciascuno dopo lo scarto dei primi 30~s. Temperatura della stanza 27\,\celsius.
Totale: 7554 campioni con presenza reale.

\paragraph{Risultato.} I valori misurati sono riportati nella
\cref{tab:fermo-seduto}.

\begin{table*}[htbp]
\centering
\caption{Scenario ``persona ferma seduta a 2,30~m''. Falsi negativi sui 7554
campioni con presenza reale, su cinque trial indipendenti da 302~s.}
\label{tab:fermo-seduto}
\small
\begin{tabular}{@{}lrr@{}}
\toprule
\textbf{Metrica} & \textbf{PIR \pir} & \textbf{Radar \ldb} \\
\midrule
Falsi negativi (media $\pm$ dev.\ std.\ su 5 trial) & $99{,}78 \pm 0{,}49\,\%$ &
  $\mathbf{0{,}00 \pm 0{,}00\,\%}$ \\
Trial con falsi negativi al 100\,\% & 4 su 5 & 0 su 5 \\
Campioni con presenza rilevata & $\sim$17 su 7554 & 7554 su 7554 \\
\bottomrule
\end{tabular}
\end{table*}

Il dato va letto con attenzione al suo significato operativo: in quattro trial su
cinque il PIR ha prodotto \textbf{zero rilevamenti in cinque minuti} con una
persona viva seduta a 2,3~m di distanza, in campo visivo libero. Il radar ha
rilevato la presenza nel 100\,\% dei campioni in tutti e cinque i trial.

Questo è il risultato che risponde alla domanda di ricerca principale. Non è un
esito sorprendente --- è precisamente ciò che la fisica del piroelettrico prevede
(\cref{sec:pir-principio}) --- ma la tesi lo \emph{quantifica}, con deviazione
standard su cinque ripetizioni, nella condizione che interessa il progetto UPRISE.

Un risultato analogo, ottenuto in una sessione precedente con procedura meno
controllata, conferma il quadro: in un tratto di \textbf{93,2~s di immobilità
continuativa} a circa 2,9~m, il radar ha rilevato la presenza nel 100\,\% dei
campioni e il PIR nello 0,2\,\% (un solo campione su 466).

\subsection{Come va interpretato il valore percentuale}
\label{sec:interpretazione-pir}

Il numero ``99,78\,\% di falsi negativi'' richiede una precisazione onesta, perché non
è una proprietà del solo sensore: dipende da due parametri sotto il nostro controllo,
il trimmer del tempo di ritenuta (tenuto al minimo, \cref{sec:hcsr501}) e il
ponticello di ritrigger (in posizione L in questa serie, \cref{sec:jumper}).

\paragraph{Che cosa misura davvero la percentuale.} In posizione L l'uscita del PIR è
un \textbf{monostabile a durata fissa}: su tutte le sessioni acquisite in questa
configurazione sono stati contati \textbf{195 impulsi con media $3{,}46 \pm 0{,}09$~s,
nessuno oltre 5~s}, indipendentemente da quanto a lungo la persona si muovesse. La
percentuale di campioni rilevati è quindi il prodotto fra il \emph{numero di eventi} e
il tempo di ritenuta, diviso la durata del trial: è un conteggio di eventi travestito
da frazione di tempo. Con la ritenuta regolata al massimo il PIR
\emph{sembrerebbe} rilevare molto di più, ma starebbe soltanto trattenendo l'ultimo
evento più a lungo.

In posizione H la percentuale cambia natura --- gli impulsi si concatenano e diventa
una vera frazione di tempo con movimento rilevato --- ma resta legata alla ritenuta,
che determina quanto a lungo un movimento intermittente venga ``riempito''. In nessuna
delle due configurazioni la percentuale è una proprietà robusta del sensore isolato.

\paragraph{Il dato robusto è il numero di eventi.} Per questo motivo, in tutto il
capitolo le percentuali sono riportate insieme alla struttura degli impulsi che le
produce. Il dato non contestabile è quello: nei 93~s di immobilità continuativa della
sessione pilota e nei cinque trial da cinque minuti di questa serie, \textbf{il PIR
non ha prodotto alcun evento di movimento} salvo le rarissime eccezioni contate sopra,
perché non c'era alcun transito di flusso infrarosso da rilevare. E un evento che non
esiste non può essere né allungato dalla ritenuta né concatenato dal ritrigger: è
questo che rende il risultato indipendente da entrambi i parametri, come verificato
per misura in \cref{tab:jumper}.

\section{Il ponticello di ritrigger: una correzione metodologica}
\label{sec:jumper}

Una parte della campagna di fase~1 è stata acquisita con l'\pir\ in una
configurazione non ottimale, e la scoperta di questo fatto ha richiesto di ripetere
alcune serie. Il problema, il modo in cui è stato individuato e la delimitazione di
ciò che ne risente sono riportati qui per intero: sono parte del risultato, non un
incidente da nascondere.

\paragraph{Il problema.} L'\pir\ dispone di un ponticello hardware che abilita o
disabilita il \emph{ritrigger} (\cref{sec:hcsr501}). In posizione L l'uscita è un
monostabile non ritriggerabile: dopo un evento resta alta per il tempo di ritenuta
impostato dal trimmer e poi ricade, \emph{ignorando} gli eventi che arrivano nel
frattempo. In posizione H ogni nuovo evento ricarica il monostabile, e con movimento
continuo l'uscita resta alta senza interruzioni. Tutte le acquisizioni fino al
22/08/2026 sono state svolte in posizione L.

\paragraph{Come è stato individuato.} Non da un'ispezione della scheda, ma dai dati.
L'analisi della durata degli impulsi sulle sessioni di fase~1 ha prodotto
\textbf{195 impulsi con media $3{,}46 \pm 0{,}09$~s, tutti compresi fra 3,40 e
3,60~s e \emph{nessuno} oltre 5~s} --- nemmeno nei trial in cui il soggetto
attraversava ripetutamente il campo per cinque minuti. Una dispersione di 0,09~s su
195 impulsi non è compatibile con un sensore che ricarica il monostabile: è la firma
di una durata fissa.

La conferma è stata diretta. Spostato il ponticello su H e ripetuta la prova della
mano, il PIR ha prodotto \textbf{un unico impulso continuo di almeno 9,2~s}, ancora
alto alla fine dell'acquisizione. Nella successiva serie di cammino a 1~m si sono
osservati impulsi singoli di almeno 61,6~s, cioè estesi all'intero trial.

\paragraph{Due errori di procedura, dichiarati.} Vanno riportati perché hanno
ritardato la diagnosi e perché il secondo è un errore di analisi, non di montaggio.

\begin{enumerate}[leftmargin=1.8em,itemsep=3pt]
  \item Due tentativi precedenti di verificare il ponticello avevano confrontato fra
        loro \emph{due posizioni entrambe diverse da H}. Il confronto non mostrava
        differenze e aveva portato alla conclusione, sbagliata, che il ponticello
        fosse inerte.
  \item La routine di estrazione degli impulsi scartava gli impulsi \emph{troncati}
        dai bordi della finestra di analisi --- cioè esattamente quelli lunghi. Sul
        file che conteneva la prova del ritrigger riportava ``nessun impulso''. La
        routine è stata corretta: gli impulsi troncati sono ora riportati a parte,
        con il loro limite inferiore di durata.
\end{enumerate}

\paragraph{Che cosa ne risente, e che cosa no.} Il ponticello agisce
\emph{solo dopo} un trigger: non può allungare un impulso che non esiste, né uno che
non riceve un secondo trigger entro il tempo di ritenuta. Ne segue che gli scenari in
cui il PIR ha prodotto zero eventi, o eventi isolati e distanti fra loro, non possono
esserne influenzati. Questa non è però rimasta una deduzione: le serie coinvolte sono
state \textbf{ripetute in posizione H} e confrontate, con l'esito della
\cref{tab:jumper}.

\begin{table*}[htbp]
\centering
\caption{Effetto del ponticello di ritrigger sul tasso di rilevamento del PIR, a
parità di scenario. Le due righe superiori cambiano molto, le due inferiori no.}
\label{tab:jumper}
\small
\begin{tabularx}{\linewidth}{@{}lrrX@{}}
\toprule
\textbf{Scenario} & \textbf{PIR in L} & \textbf{PIR in H} & \textbf{Interpretazione} \\
\midrule
Cammino sul posto, 1~m & $19{,}54\,\%$ & $\mathbf{85{,}20\,\%}$ &
  eventi ripetuti: il ritrigger li concatena \\
Sotto il banco, micro-movimenti & $35{,}86\,\%$ & $\mathbf{96{,}52\,\%}$ &
  idem, effetto ancora maggiore \\
Cammino sul posto, 2~m & $0{,}00\,\%$ & $0{,}00\,\%$ &
  nessun trigger: nulla da ricaricare \\
Sotto il banco, immobile & $0{,}10\,\%$ & $1{,}32\,\%$ &
  eventi isolati e rari: effetto trascurabile \\
\bottomrule
\end{tabularx}
\end{table*}

Le due righe inferiori sono la verifica sperimentale della deduzione, e sono state
acquisite proprio per averla. Da esse si estende, ora con il sostegno di una misura,
la stessa conclusione agli scenari di stanza vuota, di persona ferma a 2,3~m e di
movimento a 3, 4 e 5~m: tutti a zero eventi PIR, tutti insensibili al ponticello.

\paragraph{Conseguenze sulla presentazione dei risultati.} Sono tre.
\emph{(i)} I due risultati portanti della tesi --- persona ferma a 2,3~m
(\cref{sec:persona-ferma}) e persona ferma sotto il banco (\cref{sec:sotto-banco})
--- restano validi come acquisiti, e il secondo è comunque disponibile anche nella
versione in H. \emph{(ii)} Le due coppie di confronto centrali del capitolo sono
\textbf{interamente in posizione H}, quindi nella configurazione più favorevole al
PIR: non c'è alcun asterisco da apporre al confronto. \emph{(iii)} Ogni serie
riportata nel seguito indica la posizione del ponticello con cui è stata acquisita
(\cref{tab:sessioni}).

\paragraph{Nota operativa.} Dopo un ciclo di alimentazione l'\pir\ resta cieco per
circa 60~s (stabilizzazione del piroelettrico). Nella verifica del ponticello il
primo impulso è arrivato a $t = 33$~s benché il movimento fosse continuo dall'inizio:
non era un guasto. Le acquisizioni successive prevedono un transitorio iniziale
scartato in analisi, la cui durata dipende dallo scenario: 20~s in quelli ordinari,
40~s sotto il banco --- dove il soggetto deve anche posizionarsi --- e 120~s in quelli
di selettività spaziale, per una ragione scoperta soltanto in seguito e discussa in
\cref{sec:selettivita}. La durata usata è indicata per ogni serie, e dove il risultato
potrebbe dipenderne è riportata l'analisi di sensibilità.

\section{Il PIR con la persona in movimento}
\label{sec:pir-movimento}

\paragraph{Procedura.} Soggetto che cammina \emph{sul posto} a distanza fissa dalla
coppia di sensori --- movimento continuo, ma senza spostamento nello spazio. Cinque
distanze (1, 2, 3, 4, 5~m), cinque trial per distanza, 62~s utili ciascuno: 25 trial
in totale. Le serie a 1 e 2~m sono state ripetute in posizione H
(\cref{sec:jumper}), e sono quelle riportate qui.

\paragraph{Risultato.} Nella configurazione corretta il PIR \textbf{funziona bene, ma
solo da vicino}. A 1~m rileva il soggetto nell'$\mathbf{85{,}20 \pm 11{,}52\,\%}$ dei
campioni. A 2~m il tasso di rilevamento è \textbf{esattamente zero} su cinque trial, e
resta zero a 3, 4 e 5~m: venti trial consecutivi al $100\,\%$ di falsi negativi con un
soggetto in movimento continuo davanti al sensore. Il radar ha rilevato la presenza
nel $100\,\%$ dei campioni a tutte e cinque le distanze.

\paragraph{Il crollo non è graduale.} Fra 1 e 2 metri il PIR non degrada: passa da
``funziona'' a ``non vede nulla''. Non esiste una fascia intermedia in cui il
rilevamento sia parziale ma utilizzabile. Questa è la caratterizzazione corretta della
portata dell'esemplare in prova, misurata nella sua configurazione migliore e con il
trimmer di sensibilità a metà corsa.

\paragraph{Interpretazione.} Il comportamento è coerente con il meccanismo della lente
di Fresnel (\cref{sec:fresnel}): il piroelettrico risponde al \emph{transito} del
flusso infrarosso fra le zone della lente. A 1~m il soggetto sottende un angolo ampio
e le oscillazioni del busto bastano ad attraversare più zone; a 2~m lo stesso
movimento resta confinato entro una zona sola e non produce alcuna transizione. Il
fattore limitante non è la sensibilità termica del rivelatore, ma la geometria della
lente. È un'interpretazione coerente con i dati e con il principio di funzionamento,
non una misura: discriminarla da un semplice calo di sensibilità richiederebbe di
variare la distanza a parità di angolo sotteso, per esempio con un bersaglio termico
di area crescente.

\paragraph{Una affermazione precedente, corretta.} Sulla base dei dati acquisiti in
posizione L si era formulata la tesi che ``il PIR sbaglia anche con la persona in
movimento'', con un $80{,}5\,\%$ di falsi negativi misurato a 1~m. Quel valore era un
artefatto della configurazione, e l'affermazione \textbf{non è sostenibile} nella
forma generale: a 1~m e in posizione H il PIR è un buon rilevatore di movimento.
L'affermazione va riferita alla distanza, e in quella forma resta vera --- in modo
molto più netto --- oltre i 2~m.

\paragraph{Un controllo di comparabilità fra le due serie.} Poiché le serie in L e in
H sono state acquisite in momenti diversi, con il setup smontato e rimontato nel
frattempo, occorre verificare che il soggetto abbia riprodotto lo stesso movimento.
Il radar serve qui da testimone indipendente: l'energia media del bersaglio in
movimento vale 98,96 nella serie in L e 98,20 in quella in H, e la dispersione della
distanza 10,38~cm contro 11,02~cm. Il movimento è stato riprodotto quasi identico,
quindi la differenza sul PIR è attribuibile al ponticello. Resta uno scarto sulla
distanza media (99,48~cm contro 103,76~cm) dovuto al rimontaggio: va dichiarato, e
per questo motivo \textbf{le serie in H non vanno unite a quelle originali} nella
regressione della distanza di \cref{sec:distanza}, che resta costruita sui 25 trial
originali.

\section{L'esperimento controllato: la curva dose-risposta}
\label{sec:dose-risposta}

Le sezioni precedenti confrontano scenari che differiscono per più di una variabile:
distanza, postura, geometria. Questa sezione riporta l'esperimento che ne varia
\textbf{una sola}.

\paragraph{Procedura.} Soggetto in piedi a 1~m dalla coppia di sensori, ponticello in
posizione H, stesso montaggio e stessa postura in tutte le condizioni. Cambia soltanto
la quantità di movimento, su tre livelli: immobilità con sola respirazione spontanea;
micro-movimenti (piccoli aggiustamenti di braccia e busto, senza spostare i piedi);
cammino sul posto. Cinque trial per condizione, 202~s utili per le prime due e 62~s
per la terza. La distanza è 1~m, cioè quella in cui il PIR ha appena dimostrato di
funzionare bene: non è un confronto costruito a suo sfavore.

\begin{table*}[htbp]
\centering
\caption{Curva dose-risposta a 1~m, ponticello in H. Fra le tre righe cambia soltanto
la quantità di movimento del soggetto. Le due colonne di destra sono grandezze del
radar, riportate per il loro uso nell'indice di vitalità (\cref{cap:vitalita}).}
\label{tab:dose-risposta}
\small
\begin{tabular}{@{}lrrrrr@{}}
\toprule
\textbf{Condizione} & \textbf{PIR} & \textbf{dev.\ rel.} & \textbf{Radar} &
\textbf{Energia mov.} & \textbf{Disp.\ dist.} \\
\midrule
Immobile          & $1{,}52 \pm 1{,}03\,\%$   & $68\,\%$ & $\mathbf{100\,\%}$ & 78,1 & 20,2~cm \\
Micro-movimenti   & $51{,}60 \pm 21{,}72\,\%$ & $\mathbf{42\,\%}$ & $\mathbf{100\,\%}$ & 86,1 & 18,3~cm \\
Cammino sul posto & $85{,}20 \pm 11{,}52\,\%$ & $14\,\%$ & $\mathbf{100\,\%}$ & 98,2 & 11,0~cm \\
\bottomrule
\end{tabular}
\end{table*}

\paragraph{Il risultato.} Il PIR è \textbf{monotono nella quantità di movimento}, dal
$1{,}52\,\%$ all'$85{,}20\,\%$; il radar resta al $100\,\%$ in tutte e tre le
condizioni. Poiché fra le righe cambia una sola variabile, la conclusione è diretta:
il fallimento del PIR sulla persona ferma non è dovuto alla distanza (a 1~m funziona
bene), né alla configurazione (è in H), né alla taratura del trimmer (invariata fra le
righe). È dovuto \textbf{all'immobilità del soggetto}.

La struttura degli impulsi rende il contrasto più concreto della percentuale. Nella
condizione di cammino il PIR produce impulsi lunghi e concatenati: sei impulsi
completi di 14,1~s in media, fino a 29,8~s, più cinque impulsi troncati dai bordi
della finestra, il più lungo dei quali copre l'intero trial di 62~s. Nella condizione
di immobilità produce \textbf{due impulsi in totale su 1010~s di acquisizione}.

\paragraph{L'imprevedibilità nella zona intermedia.} Il dato più significativo per
un'applicazione salvavita non è la media, ma la dispersione. Nella condizione di
micro-movimenti la deviazione standard è di $\pm 21{,}72$ punti percentuali su una
media di $51{,}60$: due trial identici dal punto di vista dell'operatore danno
risultati molto diversi. In deviazione relativa il PIR è \emph{meno} ripetibile in
mezzo ($42\,\%$) che agli estremi ($68\,\%$ della condizione immobile si applica a una
media prossima a zero, dove la deviazione relativa perde significato). Vicino alla
propria soglia il PIR non è soltanto meno sensibile: è \textbf{imprevedibile}. Per un
sistema di soccorso una prestazione imprevedibile è peggiore di un limite noto e
dichiarato, perché non consente di interpretare l'assenza di segnale.

\paragraph{Robustezza rispetto alla scelta del transitorio.} La durata scartata a
inizio trial non influenza le conclusioni. Ripetendo l'analisi con scarti di 20, 30,
40 e 60~s, la condizione intermedia resta fra $51{,}5$ e $52{,}4\,\%$ e quella di
cammino fra $76{,}4$ e $85{,}2\,\%$, mentre l'immobile scende da $1{,}52\,\%$ a
$0{,}70\,\%$. L'ordinamento fra le tre condizioni non cambia mai, e il residuo di
attività nella condizione immobile è concentrato nei primi secondi, cioè
nell'assestamento del soggetto.

\paragraph{Il radar fornisce indicatori graduati.} Le ultime due colonne della
\cref{tab:dose-risposta} sono un risultato a sé, rilevante per l'obiettivo O6.
L'energia media del bersaglio in movimento cresce con la quantità di movimento
(78,1 $\to$ 86,1 $\to$ 98,2) e la dispersione della distanza cala
(20,2 $\to$ 18,3 $\to$ 11,0~cm), perché un eco più forte produce una stima di distanza
più stabile. Entrambe sono \textbf{grandezze continue e monotone}, misurate a
geometria costante su tre livelli di movimento. Sono la base sperimentale dell'indice
di vitalità del \cref{cap:vitalita}: distinguere ``presente e attivo'' da ``presente e
immobile'' non richiede un sensore diverso, richiede di leggere questi due canali
invece del solo bit di presenza.

\section{Accuratezza della misura di distanza}
\label{sec:distanza}

Gli stessi 25 trial forniscono la caratterizzazione della misura di distanza del
radar, riportata nella \cref{tab:distanza}.

\begin{table*}[htbp]
\centering
\caption{Accuratezza della distanza riportata dal \ldb, su 25 trial (5 distanze
$\times$ 5 ripetizioni, 62~s utili ciascuno). ``Misurata'' è la media fra trial con la
sua deviazione standard; ``residuo'' è lo scarto dalla retta di regressione
dell'\cref{eq:regressione}; ``disp.\ entro trial'' è la deviazione standard della
distanza all'interno di un singolo trial, mediata sui cinque. L'energia è il valore
medio del bersaglio in movimento, normalizzato 0--100.}
\label{tab:distanza}
\small
\begin{tabular}{@{}rrrrrrr@{}}
\toprule
\textbf{Reale} & \textbf{Misurata} & \textbf{Errore} & \textbf{Err.\ rel.} &
\textbf{Residuo} & \textbf{Disp.\ entro trial} & \textbf{Energia} \\
\midrule
1~m & $99{,}48 \pm 2{,}61$~cm  & $-0{,}52$~cm & $-0{,}5\,\%$ & $-3{,}0$~cm & $10{,}4$~cm & 99,0 \\
2~m & $211{,}22 \pm 0{,}77$~cm & $+11{,}22$~cm & $+5{,}6\,\%$ & $+4{,}9$~cm & $14{,}3$~cm & 85,1 \\
3~m & $309{,}76 \pm 1{,}13$~cm & $+9{,}76$~cm & $+3{,}3\,\%$ & $-0{,}4$~cm & $17{,}5$~cm & 54,4 \\
4~m & $411{,}84 \pm 2{,}76$~cm & $+11{,}84$~cm & $+3{,}0\,\%$ & $-2{,}1$~cm & $26{,}4$~cm & 35,1 \\
5~m & $518{,}20 \pm 2{,}13$~cm & $+18{,}20$~cm & $+3{,}6\,\%$ & $+0{,}5$~cm & $11{,}7$~cm & 27,6 \\
\bottomrule
\end{tabular}
\end{table*}

Due colonne meritano attenzione prima di passare al modello. L'\textbf{errore
relativo non è monotono}: il valore più alto non è a 5~m ma a 2~m ($+5{,}6\,\%$), il
che esclude la lettura ingenua di ``errore che peggiora con la distanza''. La
regressione della sezione seguente rende conto di questa irregolarità, perché tutti e
cinque i punti giacciono sulla stessa retta entro 4,9~cm.

La \textbf{dispersione entro trial}, invece, non è spiegata dal modello: cresce da
10,4 a 26,4~cm fra 1 e 4~m e poi ricade a 11,7~cm a 5~m. È una grandezza diversa
dall'errore --- misura quanto la lettura oscilla \emph{durante} un trial, non quanto è
sbagliata --- e riflette in larga parte l'oscillazione reale del busto di chi cammina
sul posto. Il valore basso a 5~m resta un'\textbf{osservazione aperta}: si può
ipotizzare che a energia ridotta (27,6, prossima alla soglia del gate corrispondente)
il modulo riporti soltanto gli echi più stabili, ma con questi dati è una congettura e
va segnalata come tale.

\subsection{Il modello di errore}

La regressione lineare fra distanza misurata e distanza reale, sui 25 trial, dà:

\begin{equation}
d_{\text{misurata}} = 1{,}0381 \cdot d_{\text{reale}} - 1{,}32~\text{cm},
\qquad R^2 = 0{,}99965
\label{eq:regressione}
\end{equation}

Il risultato si legge in due parti, ed è la lettura --- non il numero grezzo --- il
contributo di questa sezione.

\begin{itemize}[leftmargin=1.4em,itemsep=3pt]
  \item Il radar è \textbf{estremamente lineare}: $R^2 = 0{,}99965$ e residui dalla
        retta non superiori a 4,9~cm su tutto il range. Non c'è alcuna non linearità
        apprezzabile fra 1 e 5~m.
  \item L'errore è quasi interamente un \textbf{errore di scala del $+3{,}81\,\%$},
        con offset praticamente nullo ($-1{,}32$~cm). Questo spiega perché l'errore
        grezzo cresce con la distanza (da $-0{,}5$~cm a 1~m fino a $+18{,}2$~cm a
        5~m) ed è, soprattutto, \textbf{correggibile con un solo coefficiente}:
        dividendo la lettura per 1,0381 l'errore residuo scende sotto i 5~cm a tutte
        le distanze.
\end{itemize}

\paragraph{Limite dichiarato sulla causa.} I dati raccolti \emph{non} permettono di
attribuire l'origine del $+3{,}81\,\%$. Può trattarsi della calibrazione interna del
modulo oppure del riferimento con cui sono state segnate le distanze sul pavimento.
Per discriminare servirebbe un riferimento di distanza indipendente, per esempio un
distanziometro laser. La distinzione è rilevante: nel primo caso il coefficiente
correttivo è una proprietà del modulo e vale in generale, nel secondo è un artefatto
del nostro setup.

\subsection{Dispersione entro il trial}

La dispersione della distanza \emph{entro} un singolo trial è risultata di
10--26~cm nello scenario di cammino sul posto, contro
$\mathbf{1{,}48 \pm 0{,}77}$~\textbf{cm} nello scenario da seduto immobile
(\cref{sec:persona-ferma}). Il confronto è informativo: la dispersione maggiore non
è rumore del sensore, è l'oscillazione reale del busto di chi cammina sul posto. Il
radar la sta misurando correttamente.

Nello scenario da seduto emerge una scomposizione utile: la variabilità \emph{fra}
trial della distanza media è di $\pm 2{,}38$~cm (media $229{,}78$~cm su una distanza
reale di 230~cm), mentre la stabilità \emph{entro} il trial è di $1{,}48$~cm. La
variabilità di come il soggetto si siede ($\sim$2,4~cm) è quindi maggiore di quella
del sensore ($\sim$1,5~cm): il fattore limitante dell'accuratezza in questo scenario
è il soggetto, non lo strumento.

\subsection{Decadimento dell'energia con la distanza}

L'energia media del bersaglio decade da 99 a 1~m fino a 28 a 5~m
(\cref{tab:distanza}). Il decadimento è \emph{molto} più lento di quanto
prevederebbe l'equazione del radar, che per un bersaglio puntiforme dà una
dipendenza $1/D^4$ della potenza ricevuta.

Va quindi evitata una tentazione interpretativa: \textbf{questo valore non è una
potenza e non va letto con l'equazione del radar}. È un indicatore normalizzato
0--100 prodotto da un'elaborazione interna non documentata, con ogni probabilità
compressiva. Se ne può usare la monotonia (più lontano $\to$ meno energia) ma non il
valore assoluto in termini fisici.

L'osservazione ha però un uso pratico: l'extrapolazione della tendenza porta a un
valore di circa 27 a 6~m, contro una soglia di rilevamento di 15 sul gate
corrispondente (\cref{tab:soglie}). Questo indica che il sensore avrebbe funzionato
anche a 6~m.

\paragraph{Perimetro sperimentale.} Il range coperto è \textbf{1--5~m} e non 6~m: è
un limite della stanza disponibile, non del sensore, e va dichiarato come tale. La
retta di regressione è costruita su cinque punti e 25 trial. Per lo scenario UPRISE
la limitazione è comunque irrilevante, perché la distanza di interesse è
\emph{inferiore al metro} --- la persona sotto il banco --- che è coperta dallo
scenario di \cref{sec:sotto-banco}.

\section{Latenze di rilevamento e di rilascio}
\label{sec:latenza}

\paragraph{Metodo.} La ground truth registrata nei file è statica, quindi la latenza
non può essere ricavata da essa: serve un \textbf{evento a istante noto}. La
convenzione adottata è che l'acquisizione parta con la stanza vuota e che il soggetto
entri (o esca) a $t = 30$~s, con l'istante annunciato a voce dallo script di
acquisizione. L'annuncio vocale è emesso dal calcolatore che sta registrando, e conta
dalla prima riga di dati ricevuta --- cioè dalla stessa origine dei tempi che usa lo
script di analisi. Un cronometro esterno non sarebbe adeguato: fra l'apertura della
porta seriale e la prima riga passano 2--3~s fra riavvio del microcontrollore e
inizializzazione, un errore sistematico maggiore della grandezza da misurare.

Lo script di analisi riceve l'istante dell'evento come parametro e calcola la latenza
come distanza temporale fra l'evento e la prima transizione stabile di ciascun
sensore. Riporta inoltre, per ogni trial, la frazione di campioni in cui il sensore
era già attivo \emph{prima} dell'evento: un trial con questa frazione diversa da zero
è contaminato e va scartato.

\subsection{Latenza di rilevamento all'ingresso}

Dieci trial da 62~s, soggetto fuori dalla stanza all'avvio, ingresso annunciato a
$t = 30$~s. Tutti e dieci sono risultati validi: la frazione di attivazione
pre-evento è $0{,}00\,\%$ per entrambi i sensori in tutti i trial.

\begin{table*}[htbp]
\centering
\caption{Latenza di rilevamento all'ingresso, 10 trial. La terza riga è la differenza
calcolata trial per trial, non la differenza delle prime due.}
\label{tab:latenza-ingresso}
\small
\begin{tabularx}{\linewidth}{@{}lrX@{}}
\toprule
\textbf{Grandezza} & \textbf{Valore} & \textbf{Nota} \\
\midrule
Latenza radar \ldb & $5{,}36 \pm 0{,}30$~s & latenza operativa, non del solo sensore \\
Latenza PIR \pir & $5{,}96 \pm 0{,}31$~s & idem \\
Differenza appaiata radar $-$ PIR & $\mathbf{-0{,}60 \pm 0{,}35}$~\textbf{s} &
  il radar rileva per primo in 10 trial su 10 \\
\bottomrule
\end{tabularx}
\end{table*}

\paragraph{Le latenze assolute non sono latenze di sensore.} Va dichiarato subito,
perché è la differenza fra un numero corretto e un numero pubblicabile: i valori
intorno a 5--6~s sono dominati dal tragitto di rientro del soggetto e dal suo tempo di
reazione all'annuncio, non dal tempo di risposta dei sensori. Vanno riportati come
\textbf{latenza operativa}, cioè come limite superiore della latenza del sensore, e
non vanno confrontati con i valori di datasheet.

\paragraph{La grandezza pulita è la differenza appaiata.} Poiché i due sensori
osservano lo \emph{stesso} ingresso nello \emph{stesso} trial, il tragitto e il tempo
di reazione sono identici per entrambi e si cancellano nella differenza. Il valore
$-0{,}60 \pm 0{,}35$~s è quindi una proprietà dei sensori, non dell'operatore.
L'errore standard della media è $0{,}35/\sqrt{10} = 0{,}11$~s: il vantaggio del radar
vale 5,4 volte la propria incertezza. Il segno è inoltre concorde in tutti e dieci i
trial, il che sotto l'ipotesi nulla di equivalenza ha probabilità
$2 \cdot 2^{-10} \approx 0{,}002$.

\paragraph{Un limite di risoluzione da dichiarare.} Il campionamento è a 5~Hz, quindi
le latenze sono quantizzate a multipli di 200~ms e le differenze osservate valgono da
uno a sette campioni ($-0{,}2$~s in due trial, $-1{,}4$~s in uno). Nessun singolo
trial dimostrerebbe alcunché: è la \textbf{concordanza del segno} su dieci ripetizioni
a sostenere la conclusione, non l'ampiezza di una singola misura.

\paragraph{Una previsione smentita.} Ci si attendeva che in questo scenario il PIR
andasse alla pari o meglio del radar, perché attraversare una porta è esattamente lo
stimolo trasversale per cui la lente di Fresnel è progettata (\cref{sec:fresnel}).
Non è andata così. Una spiegazione plausibile è che il radar, avendo circa 6~m di
portata, agganci il soggetto mentre è ancora in avvicinamento, mentre il PIR richiede
che il flusso infrarosso attraversi le zone della lente, cosa che avviene più tardi e
più vicino. Questa spiegazione \emph{non è dimostrata} dai dati raccolti:
verificarla richiederebbe di variare la distanza della porta dal sensore.

\subsection{Latenza di rilascio all'uscita}

Cinque trial da 102~s, soggetto presente all'avvio, uscita annunciata a $t = 30$~s.
Il rilascio è stato misurato in tutti e cinque i trial per entrambi i sensori.

\begin{table*}[htbp]
\centering
\caption{Tempo di rilascio all'uscita, 5 trial. Le ``riaccensioni'' sono le
transizioni verso lo stato di presenza osservate dopo il rilascio, entro la fine del
trial.}
\label{tab:latenza-uscita}
\small
\begin{tabularx}{\linewidth}{@{}lrrX@{}}
\toprule
\textbf{Sensore} & \textbf{Rilascio} & \textbf{Riaccensioni} & \textbf{Nota} \\
\midrule
Radar \ldb & $\mathbf{18{,}36 \pm 0{,}54}$~\textbf{s} & 0 & rilascio netto in tutti i trial \\
PIR \pir & $12{,}96 \pm 1{,}40$~s & 0 & idem \\
\bottomrule
\end{tabularx}
\end{table*}

\paragraph{Qui il radar è nettamente più lento.} Impiega $5{,}40$~s in più del PIR a
dichiarare la stanza vuota. Propagando gli errori standard delle due medie, la
differenza vale $5{,}40 \pm 0{,}67$~s, cioè circa 8 volte la propria incertezza già
con cinque trial soltanto --- mentre all'ingresso ne erano serviti dieci perché
l'effetto era di $0{,}60$~s contro una dispersione di $0{,}35$. Il confronto fra i due
scenari illustra un punto metodologico generale: il numero di ripetizioni necessarie
non dipende dallo scenario ma dal rapporto fra l'effetto e la sua variabilità.

\paragraph{Nessuna riaccensione.} In tutti e cinque i trial entrambi i sensori
rilasciano una sola volta e restano bassi. È un comportamento migliore di quello
osservato nella sessione di validazione (\cref{sec:coda}), dove dopo l'uscita il radar
aveva mantenuto la presenza per circa 10~s riportando un bersaglio fermo fantasma a
5,2~m con energia in decadimento.

\paragraph{Stima della coda propria del radar.} Il tempo di rilascio misurato include
il tempo che il soggetto impiega a uscire dal campo, che non è noto direttamente. Il
PIR può però essere usato come cronometro: la sua uscita ricade $3{,}46 \pm 0{,}09$~s
dopo l'ultimo trigger (\cref{sec:jumper}), quindi l'istante di uscita dal campo è
circa $12{,}96 - 3{,}5 = 9{,}5$~s dopo l'annuncio, e la coda propria del radar è circa
$18{,}36 - 9{,}5 \approx \mathbf{8{,}9}$~\textbf{s}.

Questa è una \emph{stima}, non una misura, e il suo limite va dichiarato: i due
sensori hanno campi di vista diversi, quindi i rispettivi istanti di ``uscita dal
campo'' non coincidono esattamente. Il valore è però coerente con le altre due
osservazioni indipendenti disponibili: i circa 10~s della sessione di validazione e i
circa 12~s di un trial preliminare.

\paragraph{Il parametro dichiarato non descrive il comportamento osservato.} Il
timeout di presenza letto dal modulo via UART vale \textbf{5~s}
(\cref{tab:identita}), mentre tre osservazioni indipendenti danno fra 9 e 12~s. Non è
un difetto --- il ritardo è una scelta di progetto che evita il lampeggio dell'uscita
--- ma il valore configurato non è utilizzabile per prevedere il comportamento. Per
lo scenario UPRISE la conseguenza è concreta: una stanza appena liberata continua a
risultare occupata per quasi dieci secondi, e un sistema che conti gli occupanti
aggregando lo stato di più sensori deve tenerne conto.

\section{Lo scenario UPRISE: persona sotto il banco}
\label{sec:sotto-banco}

È lo scenario che replica la condizione d'interesse del progetto: soggetto
accovacciato sotto un piano, a distanza inferiore al metro dal sensore.

\paragraph{Procedura.} Coppia di sensori fissata sotto il piano di un tavolo, soggetto
rannicchiato sotto di esso a circa 50--60~cm. Cinque trial da 302~s utili per
condizione, scartando 40~s di transitorio iniziale --- più lunghi che negli altri
scenari perché il soggetto deve anche posizionarsi. Due condizioni: immobilità con
sola respirazione spontanea, e presenza di micro-movimenti. Entrambe acquisite con il
ponticello del PIR in posizione H (\cref{sec:jumper}), quindi \textbf{perfettamente
appaiate} e nella configurazione più favorevole al PIR.

\begin{table*}[htbp]
\centering
\caption{Persona sotto il banco a circa 60~cm, ponticello in H. Fra le due righe
cambia soltanto la presenza di micro-movimenti.}
\label{tab:sotto-banco}
\small
\begin{tabular}{@{}lrrr@{}}
\toprule
\textbf{Condizione} & \textbf{Rilevamento PIR} & \textbf{FN PIR} & \textbf{FN radar} \\
\midrule
Con micro-movimenti & $96{,}52 \pm 3{,}60\,\%$ & $3{,}48\,\%$ & $\mathbf{0{,}00 \pm 0{,}00\,\%}$ \\
Immobile & $1{,}32 \pm 0{,}86\,\%$ & $\mathbf{98{,}68\,\%}$ & $\mathbf{0{,}00 \pm 0{,}00\,\%}$ \\
\bottomrule
\end{tabular}
\end{table*}

\paragraph{Il risultato.} Il contrasto è ancora più netto che a 1~m: $96{,}5\,\%$
contro $1{,}3\,\%$. A 60~cm il PIR è un rilevatore di movimento quasi perfetto --- e
resta \textbf{completamente cieco} alla persona ferma. Il radar rileva la presenza nel
$100\,\%$ dei campioni in entrambe le condizioni, su 8054 campioni per condizione.

Questo è il risultato che giustifica la tesi, perché è misurato nella geometria reale
del progetto e non in una configurazione di laboratorio. Insieme alla curva
dose-risposta di \cref{sec:dose-risposta} costituisce la \textbf{seconda coppia
appaiata} del capitolo: due distanze diverse (1~m e 0,6~m), due geometrie diverse (in
piedi e accovacciato sotto un piano), due volte la stessa conclusione.

\paragraph{La versione in posizione L, per completezza.} La stessa coppia era stata
acquisita in precedenza con il ponticello in L, dando $35{,}86\,\%$ con
micro-movimenti e $0{,}10\,\%$ da immobile. Il ritrigger cambia molto la prima riga e
quasi nulla la seconda, come previsto e come verificato in \cref{tab:jumper}. Le due
serie non vanno però mescolate: fra l'una e l'altra il montaggio è stato rifatto, e la
distanza media riportata dal radar passa da $62{,}9 \pm 4{,}2$~cm a
$71{,}1 \pm 4{,}6$~cm. Lo scarto di circa 8~cm è del setup, non del sensore.

\subsection{Il radar rileva bene ma localizza male}

La dispersione della distanza \emph{entro} il trial vale 14,4~cm nella serie in L e
24,2~cm in quella in H, contro i $1{,}48 \pm 0{,}77$~cm misurati sul soggetto seduto
immobile a 2,3~m (\cref{sec:persona-ferma}): circa \textbf{dieci volte peggio}, sullo
stesso tipo di bersaglio immobile e a distanza minore.

Sotto l'arredo il radar rileva quindi la presenza in modo impeccabile ma stima la
distanza in modo molto più rumoroso. Per lo scenario UPRISE la cosa è irrilevante ---
al soccorritore serve sapere \emph{se} c'è qualcuno sotto quel banco, non a quanti
centimetri --- ma va dichiarata, perché esclude l'uso della distanza riportata come
misura fine in questa geometria.

\subsection{La distribuzione per-gate sotto il piano}
\label{sec:gate-sotto-banco}

L'energia per-gate in questa geometria merita un commento, perché a una prima lettura
sembra incoerente con la distanza riportata. I valori medi sui cinque trial sono
riportati nella \cref{tab:gate-sotto-banco}.

\begin{table*}[htbp]
\centering
\caption{Energia media per-gate con il soggetto immobile sotto il banco a circa
60~cm, confrontata con il soggetto seduto immobile a 2,3~m. Ogni gate copre 75~cm.}
\label{tab:gate-sotto-banco}
\small
\begin{tabular}{@{}llrrrrrrrrr@{}}
\toprule
\textbf{Scenario} & \textbf{Canale} & \textbf{g0} & \textbf{g1} & \textbf{g2} &
\textbf{g3} & \textbf{g4} & \textbf{g5} & \textbf{g6} & \textbf{g7} & \textbf{g8} \\
\midrule
Sotto il banco & stazionario & 0,0 & 0,0 & 100,0 & 93,3 & 68,8 & 34,3 & 32,3 & 27,4 & 23,8 \\
$\sim$60~cm    & movimento   & 68,9 & 88,1 & 50,7 & 12,6 & 7,8 & 5,7 & 5,5 & 5,2 & 5,1 \\
\midrule
Seduto         & stazionario & 0,0 & 0,0 & 90,6 & 100,0 & 99,1 & 18,2 & 10,5 & 8,4 & 6,5 \\
2,3~m          & movimento   & 17,1 & 13,3 & 8,9 & 16,0 & 10,1 & 4,7 & 4,6 & 4,7 & 2,4 \\
\bottomrule
\end{tabular}
\end{table*}

Il quadro si legge in tre punti, e corregge un'osservazione formulata prima di
disporre di questi valori.

\begin{enumerate}[leftmargin=1.8em,itemsep=3pt]
  \item Il canale \textbf{in movimento} è del tutto coerente con la distanza: il picco
        è sui gate 0--1, cioè entro 1,5~m, come deve essere per un bersaglio a 60~cm.
        Nel confronto, sul soggetto a 2,3~m gli stessi gate valgono 17,1 e 13,3.
  \item Il canale \textbf{stazionario} non è anomalo, è \emph{troncato}: i gate 0 e 1
        sono strutturalmente nulli (\cref{sec:senergy-zero}), quindi il gate 2 è
        semplicemente il primo gate disponibile e la saturazione a 100 vi si accumula.
        Non c'è alcun bersaglio a 1,5--2,25~m da spiegare.
  \item La \textbf{coda} verso i gate lontani è invece il fatto notevole: sotto il
        banco l'energia stazionaria vale ancora 23,8 al gate 8 (cioè intorno ai 6~m),
        contro 6,5 nello scenario seduto. L'eco molto forte del bersaglio ravvicinato
        si distribuisce su tutte le celle di distanza, verosimilmente per cammini
        multipli sotto il piano. È una spiegazione plausibile e coerente con il
        confronto fra le due righe, \textbf{non una misura}: distinguerla da un
        artefatto dell'elaborazione interna richiederebbe l'accesso al segnale grezzo,
        che il modulo non espone.
\end{enumerate}

Il punto operativo è comunque acquisito, ed era una delle domande aperte dello
scenario: sotto il metro il respiro e l'indice di vitalità devono usare i
\textbf{canali in movimento}, che qui sono ricchi e correttamente collocati, non
quelli stazionari.

\subsection{Il respiro nella geometria del progetto}

Alla stessa serie è stata applicata l'analisi del respiro descritta in
\cref{sec:respiro}. Con il criterio di accettazione adottato (rapporto
segnale-rumore $> 3$, tolleranza di concordanza del $10\,\%$, almeno tre canali in
movimento concordi) il respiro è risultato estraibile in \textbf{quattro trial su
cinque}, con stima aggregata di $\mathbf{21{,}0 \pm 2{,}5}$ \textbf{atti al minuto}.
Il valore è più alto e più disperso di quello misurato sul soggetto seduto (18--20
atti/min), il che è plausibile per una postura rannicchiata.

\paragraph{Analisi di sensibilità al criterio.} Il trial escluso merita un
approfondimento, perché la sua esclusione dipende da una soglia che abbiamo scelto
noi. Abbassando il rapporto segnale-rumore minimo da 3 a 2,5 si recuperano
\textbf{cinque trial su cinque}, e --- punto essenziale --- le stime dei quattro trial
già validi restano \textbf{identiche} (17,9 / 20,9 / 23,9 / 21,4 atti al minuto). A
cambiare è soltanto l'aggregato, che diventa $21{,}9 \pm 2{,}9$ invece di
$21{,}0 \pm 2{,}5$. Il segnale nel trial escluso c'era: a scartarlo era la soglia, non
il sensore.

Questo è il modo corretto di riportare un risultato che dipende da un parametro
arbitrario: si dichiara il criterio usato, si mostra l'effetto di variarlo, e si
verifica che le conclusioni non ne dipendano. Va però ribadito il limite già indicato
in \cref{sec:respiro}: in assenza di ground truth queste restano stime concordi fra
canali, non misure con un errore. La prova valida è quella a ritmo imposto
(\cref{sec:respiro-metronomo}).

\subsection{La questione della lamiera}

Il piano del banco UPRISE è rinforzato con lamiera forata antisfondamento, e il
metallo è il materiale peggiore per un radar (\cref{sec:ostacoli}). La prova di
penetrazione attraverso la lamiera \textbf{non è stata inclusa in questo scenario}, e
la ragione è geometrica: nel progetto il sensore è montato \emph{sotto} il piano e la
persona si rifugia anch'essa sotto il piano, quindi fra sensore e soggetto non c'è
metallo. La lamiera sta sopra entrambi.

Questo non chiude la questione, la sposta: resta aperto se la lamiera al di sopra
introduca riflessioni che alterino la misura, ed è una domanda distinta da quella
della penetrazione. La coda per-gate osservata in \cref{sec:gate-sotto-banco} rende la
domanda concreta. Va posta al relatore insieme alla scelta del punto di montaggio, ed
è trattata come questione di progetto nel \cref{cap:conclusioni}.

\section{Quante persone vede il radar}
\label{sec:due-persone}

Il \ldb\ non dispone di un array di antenne e non fornisce informazione angolare. Da
questo si deduce comunemente che riporti un solo bersaglio. La formulazione è
imprecisa, e questa sezione la sostituisce con una misura.

\paragraph{Il meccanismo da verificare.} Il modulo espone \textbf{due canali
indipendenti}, uno per il bersaglio in movimento e uno per quello stazionario, ciascuno
con la propria distanza e la propria energia (\cref{sec:protocollo-seriale}). Se i due
canali possono agganciare bersagli diversi, il modulo separa due persone \emph{in
stati diversi} pur senza risoluzione angolare. L'ipotesi è verificabile.

\paragraph{Procedura.} Persona A ferma a 2~m, persona B a 4~m, tre trial da 102~s utili
per geometria. Le due persone sono \textbf{sfalsate lateralmente di circa 50~cm}
(14\textdegree\ fuori asse a 2~m, 7\textdegree\ a 4~m: entrambe ben dentro i
$\pm 60$\textdegree\ del diagramma di rilevamento del manuale). Lo sfalsamento è
necessario: in fila, il corpo davanti fa da schermo e il risultato non distinguerebbe
``il modulo non separa i bersagli'' da ``il secondo bersaglio era in ombra''. La
variante in fila è stata poi acquisita apposta, come confronto.

\paragraph{Un controllo preliminare indispensabile.} Prima di attribuire al modulo
l'incapacità di riportare B, va escluso che il problema sia la distanza. Un trial di
controllo con \textbf{B da sola, ferma a 4~m} dà rilevamento al $100\,\%$, distanza
$388{,}9 \pm 3{,}7$~cm ed energia stazionaria 99,3. Il modulo vede benissimo una
persona ferma a 4~m: il confondente è escluso.

\begin{table*}[htbp]
\centering
\caption{Capacità di separare due persone, su tre geometrie. ``B riportata'' è la
frazione di campioni in cui la persona lontana compare su almeno uno dei due canali,
riconosciuta dalla banda di distanza. Tre trial per geometria, 1833 campioni utili
ciascuna.}
\label{tab:due-persone}
\small
\begin{tabular}{@{}lrrrrr@{}}
\toprule
\textbf{Geometria} & \textbf{Entrambi i} & \textbf{B} & \textbf{Dist.} &
\textbf{Dist.} & \textbf{Presenza} \\
 & \textbf{canali attivi} & \textbf{riportata} & \textbf{staz.} & \textbf{mov.} &
\textbf{radar} \\
\midrule
B cammina, sfalsate di lato & $79{,}2\,\%$ & $\mathbf{73{,}8\,\%}$ & 212~cm & 389~cm & $100\,\%$ \\
Entrambe ferme, sfalsate    & $21{,}2\,\%$ & $\mathbf{19{,}5\,\%}$ & 220~cm & 375~cm & $100\,\%$ \\
B cammina, in fila dietro A & $42{,}4\,\%$ & $\mathbf{0{,}3\,\%}$  & 215~cm & 215~cm & $100\,\%$ \\
\bottomrule
\end{tabular}
\end{table*}

\subsection{I due canali sono lo strumento di separazione}

\paragraph{Stati diversi: la separazione funziona.} Con A ferma a 2~m e B che cammina
a 4~m, entrambi i canali sono attivi contemporaneamente nel $79{,}2\,\%$ dei campioni,
e riportano due distanze separate da 177~cm: 212~cm sul canale stazionario (A) e
389~cm sul canale in movimento (B). L'energia media del canale in movimento vale 38,4,
coerente con il 35,1 misurato a 4~m nella caratterizzazione della distanza
(\cref{tab:distanza}): è una conferma indipendente che il canale stava inseguendo la
persona lontana e non un artefatto vicino.

Il limite del modulo va quindi riformulato: non ``riporta un solo bersaglio'', ma
\textbf{riporta un bersaglio per canale}.

\paragraph{Stesso stato: la separazione degrada ma non si annulla.} Con entrambe le
persone ferme, B compare comunque nel $19{,}5\,\%$ dei campioni --- nel $7{,}2\,\%$
sul canale stazionario e per il resto sul canale in movimento, quando i suoi
micro-movimenti involontari la registrano come bersaglio mobile. Anche questa
formulazione corregge un'affermazione precedente: dire che due persone nello stesso
stato ``non vengono separate'' è troppo assoluto.

Il meccanismo si legge nell'istogramma delle distanze stazionarie: quando i due
bersagli competono per lo stesso canale, \textbf{la più vicina vince nel $92{,}8\,\%$
dei campioni}. La dispersione lo conferma: la deviazione standard della distanza
stazionaria passa da 3,7~cm (una sola persona ferma, trial di controllo) a 40,9~cm
(due ferme). Non è rumore, sono le escursioni del canale verso il bersaglio lontano.

\paragraph{In fila: il corpo davanti nasconde completamente quello dietro.} Con B
esattamente dietro A sull'asse del sensore, i due canali riportano \emph{lo stesso}
bersaglio: la differenza mediana fra le due distanze è di 3~cm, ed è inferiore a 30~cm
nel $99\,\%$ dei campioni con entrambi i canali attivi (contro il $12\,\%$ nel caso
sfalsato). B è riportata nello $0{,}3\,\%$ dei campioni, cioè praticamente mai.
L'energia più alta (58,1 contro 38,4) e la dispersione più bassa (14,2 contro 49,1~cm)
confermano che il radar ha agganciato saldamente un unico bersaglio vicino.

\subsection{Che cosa si perde, e che cosa no}

Va sottolineato il dato che attraversa tutte e tre le righe della
\cref{tab:due-persone}: la \textbf{presenza è rilevata al $100\,\%$ in ogni
geometria}. Il radar non sbaglia mai a dire ``c'è qualcuno''. Quello che perde è il
\textbf{conteggio}, non la presenza --- e il caso peggiore non è ``due persone ferme''
($19{,}5\,\%$) ma ``una dietro l'altra'' ($0{,}3\,\%$).

\paragraph{Conseguenza per UPRISE.} Un singolo sensore non può censire più persone in
una stanza: chi sta dietro a qualcun altro è invisibile. Questo però
\textbf{rafforza} l'architettura del progetto invece di indebolirla. UPRISE prevede un
sensore per arredo, ciascuno rivolto al proprio occupante: è precisamente la
configurazione in cui il limite non si manifesta, perché ogni sensore ha un solo
bersaglio da riportare. Il limite va comunque dichiarato, perché esclude un'altra
architettura --- un sensore singolo usato come contatore di presenze in aula --- che
sarebbe più economica e che i dati mostrano non essere praticabile.

\section{Selettività spaziale: isolare il proprio banco}
\label{sec:selettivita}

Nell'architettura UPRISE ogni arredo ha il proprio sensore. Perché il sistema
funzioni, ciascun sensore deve riportare il \emph{proprio} occupante e non quello del
banco accanto: un'informazione attribuita all'arredo sbagliato indirizzerebbe i
soccorritori nel posto sbagliato. Questa sezione misura quanto la selettività
spaziale sia ottenibile con i parametri del modulo.

\paragraph{Lo strumento disponibile.} Il \ldb\ permette di ridurre il
\textbf{gate massimo}, cioè la cella di distanza più lontana considerata nel
rilevamento. La regola \emph{portata $=$ gate $\times$ 75~cm} è documentata dal
costruttore, con un esempio numerico che coincide con la nostra misura
(\cref{sec:soglie}). Impostando il gate massimo a 2 si ottiene una portata di 150~cm,
adatta a coprire lo spazio sotto un banco e non oltre.

\paragraph{Il limite noto dello strumento.} Il gate massimo limita la
\textbf{distanza}, non l'\textbf{angolo}. Un vicino laterale entro 150~cm resta in
linea di principio dentro la zona coperta. La domanda sperimentale è quindi doppia:
il limite di distanza funziona davvero, e il vicino laterale disturba?

\paragraph{Perché non si è usato il gate 1.} La persona sotto il banco è stata
misurata a 62--69~cm (\cref{sec:sotto-banco}), quindi un gate massimo pari a~1, con
portata di 75~cm, la coprirebbe escludendo un vicino a un metro: sarebbe l'unica
impostazione capace di isolare davvero un banco, dato che il gate~2 arriva già a
150~cm. Non è stata adottata perché il protocollo ufficiale \textbf{si contraddice
al proprio interno} sul valore minimo ammesso. Il \S1.2.2 (p.~5) afferma che
\emph{``the range can be set from 1 to 8''}, e il manuale~\cite{hilinkManualV104}
riporta al \S5.2 un paragrafo pressoché identico; il \S2.2.3 (p.~8), che definisce
il comando \texttt{0x0060} e la codifica del suo parametro, dichiara invece
\emph{``(configuration range 2\textasciitilde8)''}. Poiché le due formulazioni
compaiono nello stesso documento e nella stessa revisione, non sono ordinabili
cronologicamente: la sola discriminante disponibile è di natura normativa, e la
sezione che definisce il comando prevale su quella descrittiva.

Sul nostro esemplare il gate~1 risponde correttamente, con portata misurata di
75~cm, coerente con la regola \emph{portata $=$ gate $\times$ 75~cm} e con
l'esempio numerico presente in entrambe le sezioni descrittive. Questa osservazione
non è però sufficiente a sanare la discrepanza, perché il firmware
\textbf{non convalida il parametro ricevuto}: il gate~0, escluso da ogni
formulazione, viene accettato senza segnalazione di errore e produce un
comportamento anomalo --- distanza riportata costantemente a 72~cm e presenza
sempre attiva. L'accettazione del comando non costituisce quindi prova di
conformità; ciò che resta significativo è la differenza fra i due casi, con il
gate~0 accettato e scorretto e il gate~1 accettato e conforme alla regola
documentata. Il gate~1 è pertanto riportato come misura su un singolo esemplare e
non è proposto come scelta di progetto per il dispositivo del progetto UPRISE.

\paragraph{Procedura.} Gate massimo impostato a 2 e riletto dal modulo per conferma.
Quattro scenari, tre trial da 182~s ciascuno, occupante sempre \emph{fermo}:

\begin{enumerate}[leftmargin=1.8em,itemsep=2pt]
  \item occupante a 1~m sull'asse, da solo (riferimento);
  \item nessun occupante, una persona a 3~m sull'asse, cioè oltre il limite di 150~cm;
  \item nessun occupante, una persona a 1~m in direzione \emph{perpendicolare}
        all'asse del sensore;
  \item occupante a 1~m sull'asse \emph{più} il vicino perpendicolare a 1~m.
\end{enumerate}

\subsection{Il limite di distanza funziona}

Negli scenari 2 e 3 il tasso di rilevamento del radar misurato con il transitorio
standard di 20~s vale rispettivamente $0{,}00\,\%$ e $24{,}83 \pm 19{,}55\,\%$. Il
secondo valore sembrerebbe indicare che il vicino perpendicolare viene visto quasi un
quarto del tempo. \textbf{Non è così}, e la verifica merita di essere riportata perché
è un caso in cui la metrica aggregata induce in errore.

Esaminando l'andamento temporale, in tutti e tre i trial dello scenario 3 il radar
risulta attivo \emph{dall'istante zero}, si spegne una sola volta --- dopo 65,6, 29,4
e 100,8~s rispettivamente --- e \textbf{non si riaccende mai più} per i restanti
100--170~s, benché una persona viva si trovi a 1~m di distanza. Non è un rilevamento
intermittente: è una coda di presenza residua dalla fase di posizionamento, che decade
e non si ripresenta. Il numero di fronti di salita per ora è infatti $0{,}00$ in tutti
i trial.

Restringendo l'analisi allo stato stazionario (scartando i primi 120~s, oltre la coda
più lunga osservata) il quadro è netto, ed è riportato nella
\cref{tab:selettivita}: il vicino a 3~m e il vicino perpendicolare a 1~m sono
\textbf{entrambi al $0{,}00\,\%$}, mentre l'occupante sull'asse resta al $100\,\%$.

\begin{table*}[htbp]
\centering
\caption{Selettività spaziale con gate massimo 2 (portata 150~cm), tre trial per
scenario. La colonna ``a regime'' scarta i primi 120~s, oltre la coda di assestamento
più lunga osservata; è la colonna che descrive il comportamento del sensore.}
\label{tab:selettivita}
\small
\begin{tabularx}{\linewidth}{@{}lrrX@{}}
\toprule
\textbf{Scenario} & \textbf{Finestra std.} & \textbf{A regime} & \textbf{Lettura} \\
\midrule
Occupante a 1~m sull'asse & $100{,}00\,\%$ & $\mathbf{100{,}00\,\%}$ &
  rilevato sempre \\
Persona a 3~m sull'asse & $0{,}00\,\%$ & $\mathbf{0{,}00\,\%}$ &
  il limite di distanza funziona \\
Persona a 1~m a 90\textdegree & $24{,}83\,\%$ & $\mathbf{0{,}00\,\%}$ &
  coda di assestamento, non rilevamento \\
Occupante $+$ vicino a 90\textdegree & $100{,}00\,\%$ & $\mathbf{100{,}00\,\%}$ &
  vedi confronto appaiato \\
\bottomrule
\end{tabularx}
\end{table*}

\paragraph{Un'osservazione metodologica.} Il transitorio di 20~s adottato negli altri
scenari \textbf{non è sufficiente} in questo, dove la coda di presenza dalla fase di
posizionamento arriva a superare i 100~s. Il valore ``$24{,}83 \pm 19{,}55\,\%$''
calcolato con la convenzione standard è un artefatto del protocollo, non una proprietà
del sensore, e la deviazione standard enorme rispetto alla media era il primo indizio.
Negli scenari di selettività il transitorio va portato a 120~s.

\subsection{Il vicino laterale non sporca la lettura del proprio occupante}

Resta la domanda più importante: con l'occupante presente, il vicino perpendicolare
altera la misura? Il confronto è \textbf{appaiato} --- stesso soggetto, stessa
posizione, stessa distanza, cambia solo la presenza del vicino --- ed è riportato
nella \cref{tab:vicino}.

\begin{table*}[htbp]
\centering
\caption{Effetto di un vicino fermo a 1~m in direzione perpendicolare, sulla lettura
dell'occupante fermo a 1~m sull'asse. Tre trial per scenario, finestra standard. La
significatività è il rapporto fra la differenza e l'errore standard propagato.}
\label{tab:vicino}
\small
\begin{tabular}{@{}lrrr@{}}
\toprule
\textbf{Grandezza} & \textbf{Occupante da solo} & \textbf{Con vicino} &
\textbf{Significatività} \\
\midrule
Distanza stazionaria media & $102{,}37 \pm 2{,}80$~cm & $106{,}93 \pm 6{,}56$~cm & $1{,}1\,\sigma$ \\
Dispersione della distanza & $14{,}47 \pm 0{,}93$~cm & $16{,}07 \pm 1{,}31$~cm & $1{,}7\,\sigma$ \\
Distanza in movimento media & $101{,}73 \pm 2{,}81$~cm & $105{,}70 \pm 6{,}75$~cm & $0{,}9\,\sigma$ \\
Energia media in movimento & $87{,}73 \pm 2{,}19$ & $83{,}87 \pm 2{,}61$ & $2{,}0\,\sigma$ \\
Rilevamento della presenza & $100{,}00\,\%$ & $100{,}00\,\%$ & --- \\
\bottomrule
\end{tabular}
\end{table*}

\paragraph{Nessuna differenza significativa.} Con tre trial per scenario nessuna
grandezza supera le $2\,\sigma$: la perturbazione introdotta dal vicino è al più di
pochi centimetri sulla distanza e di pochi punti sull'energia. Restringendo l'analisi
allo stato stazionario le differenze si riducono ulteriormente (tutte sotto
$1{,}5\,\sigma$), quindi la conclusione non dipende dalla finestra scelta.

Va però dichiarato che tre trial danno poca potenza statistica: l'affermazione
sostenibile è ``non si osserva una perturbazione superiore a pochi centimetri'', non
``non c'è perturbazione''. Per il caso d'uso la distinzione è irrilevante --- serve
sapere \emph{se} il banco è occupato --- ma la formulazione va tenuta corretta.

\subsection{Conta l'angolo, non la distanza}

Il confronto fra questo esperimento e quello con due persone a 2 e 4~m
(\cref{sec:due-persone}) individua la variabile che governa il fenomeno.

\begin{itemize}[leftmargin=1.4em,itemsep=3pt]
  \item Con due persone \textbf{sfalsate di soli 50~cm} --- cioè 14\textdegree\ e
        7\textdegree\ fuori asse, entrambe in pieno lobo principale --- la dispersione
        della distanza stazionaria passava da 3,7 a 40,9~cm: il secondo bersaglio
        competeva eccome per il canale.
  \item Con il secondo corpo a \textbf{90\textdegree}, alla stessa distanza
        dell'occupante, la dispersione passa da 14,5 a 16,1~cm e il vicino da solo non
        è nemmeno riportato.
\end{itemize}

La differenza fra i due casi non è la distanza, che nel secondo è addirittura minore:
è l'\textbf{angolo}. Un bersaglio molto fuori asse è attenuato dal diagramma di
irradiazione e non compete con un occupante che satura il canale.

\paragraph{Due limiti da dichiarare, e che cosa resta aperto.} In entrambi gli
scenari 3 e 4 il vicino era \textbf{fermo}. Un vicino \emph{in movimento} a
90\textdegree\ è una condizione diversa e non è stata misurata: il canale in movimento
è libero, e i dati di \cref{sec:due-persone} mostrano che è proprio il canale che
aggancia più facilmente un secondo bersaglio. È la prima estensione da fare su questo
scenario.

Inoltre, durante la preparazione degli esperimenti di latenza si era osservato che il
modulo rileva un soggetto posto \emph{lateralmente e arretrato} rispetto al proprio
asse, in una posizione che il diagramma del manuale darebbe fuori copertura, mentre in
asse alle spalle non lo rilevava. L'osservazione è qualitativa e la geometria non era
misurata: il dato accertato è soltanto che la copertura eccede il lobo principale sui
lati, \textbf{non} che il modulo sia omnidirezionale. I risultati di questa sezione
non la contraddicono --- a 90\textdegree\ e con bersaglio fermo il vicino non è
riportato --- ma non la chiudono: serve una caratterizzazione angolare dedicata, con
il soggetto a distanza fissa e angolo variato per passi noti, che è pianificata e non
ancora svolta.

\paragraph{Conseguenza per UPRISE.} Con gate massimo 2 un sensore per banco riporta il
proprio occupante in modo affidabile e non è corrotto da un vicino fermo, né in asse
oltre la portata né a 90\textdegree. Il rischio residuo non è la corruzione della
lettura, è l'\textbf{attribuzione errata}: un banco vuoto il cui sensore rilevasse il
vicino segnalerebbe un occupante inesistente. Nella condizione provata questo non
avviene, ma il caso del vicino in movimento resta da verificare prima di considerare
chiusa la questione.

\section{Penetrazione degli ostacoli}
\label{sec:ostacoli}

\emph{Acquisizione in corso.}

\paragraph{Metodo previsto.} Ripetizione dello scenario ``persona ferma'' con
l'interposizione di un pannello fra sensore e soggetto, per i materiali disponibili:
cartongesso, legno, vetro, plastica. Metrica: tasso di falsi negativi e riduzione
dell'energia del bersaglio rispetto al caso senza ostacolo, a parità di distanza.

\paragraph{Attese da letteratura.} Il radar a 24~GHz attraversa cartongesso, legno,
vetro e plastica con attenuazione contenuta, mentre è bloccato da metallo e da
cemento armato di spessore rilevante. Il PIR è invece bloccato da \emph{tutti} i
materiali elencati, vetro compreso: la finestra ottica 5--14~\textmu m del corpo
umano non attraversa il vetro comune. Quest'ultimo punto è un risultato atteso ma
importante per il confronto, perché rende il PIR inutilizzabile per qualunque
montaggio incassato o protetto.

\section{Il rilevamento del respiro}
\label{sec:respiro}

\subsection{Metodo}

Il metodo segue quanto descritto in \cref{sec:respiro-teoria}: si campiona a 5~Hz
un canale di energia per-gate in engineering mode, si applica la FFT alla serie
temporale e si cerca il picco in banda 0,1--0,5~Hz. La frequenza del picco,
moltiplicata per 60, dà la stima degli atti al minuto.

Il contributo metodologico di questa tesi è nella \textbf{selezione del canale}. Il
LD2410B mette a disposizione venti serie temporali di energia (due di bersaglio più
diciotto per-gate) e non tutte sono utilizzabili: alcune sono sature
(\cref{sec:saturazione}), altre identicamente nulle (\cref{sec:senergy-zero}), altre
piatte perché lontane dal bersaglio. La modalità \texttt{-{}-scan} dello script di
analisi esamina tutti i canali, scarta quelli saturi e quelli piatti secondo criteri
dichiarati, ordina i rimanenti per rapporto segnale-rumore e riporta la
\textbf{mediana delle stime concordi}. La concordanza fra canali indipendenti è
l'unico controllo di plausibilità disponibile in assenza di ground truth.

\subsection{Prima prova: soggetto a 40~cm}

\paragraph{Risultato.} Soggetto fermo a circa 40~cm, 91~s di acquisizione,
respirazione spontanea. \textbf{Sette canali di energia in movimento, indipendenti
fra loro, concordano su 0,264~Hz, cioè 15,8 atti al minuto}, con rapporto
segnale-rumore fino a 12,8 sul canale \texttt{menergy\_gate2}. Il valore è
fisiologicamente plausibile per un adulto a riposo.

\paragraph{Perché non è una misura valida.} Va detto con chiarezza, perché è una
distinzione che qualifica il lavoro: questa prova dimostra la \emph{fattibilità} ma
non è una misura utilizzabile come risultato. Le ragioni sono tre:
\emph{(i)} manca la ground truth, cioè non si conosce il ritmo respiratorio reale e
quindi non si può calcolare un errore; \emph{(ii)} l'89\,\% dei campioni sui canali
in movimento è saturo; \emph{(iii)} 40~cm non è una distanza rappresentativa di
alcuno scenario d'uso.

\paragraph{I canali stazionari sono inutilizzabili a questa distanza.} I gate 2 e 3
stazionari sono saturi al 100\,\%, i gate 0 e 1 sono identicamente nulli
(\cref{sec:senergy-zero}) e i gate 4--8 forniscono stime fra 6,6 e 9,2 atti al
minuto, in disaccordo fra loro e con il canale in movimento. Il valore 0,132~Hz che
appare su alcuni di questi canali è esattamente la metà di 0,264~Hz: si tratta con
ogni probabilità di una subarmonica, cioè di un artefatto dell'analisi, non di un
segnale.

\subsection{Seconda prova: soggetto a 2,31~m}

La stessa analisi è stata applicata ai cinque trial dello scenario ``persona ferma''
(\cref{sec:persona-ferma}), a 2,31~m e per 302~s ciascuno. Il quadro si ripete, ma
in condizioni migliori.

\begin{itemize}[leftmargin=1.4em,itemsep=3pt]
  \item I canali \textbf{in movimento} dei gate 2--8 concordano su circa 0,30~Hz,
        cioè \textbf{18,2 atti al minuto}, valore plausibile per un adulto seduto.
        Le stime per i cinque trial sono 18,2 / 18,9 / 18,8 / 20,0 / 20,2 atti al
        minuto, con 7--8 canali concordi per trial.
  \item Questi canali \textbf{non sono saturi} a questa distanza: la percentuale di
        campioni a fondo scala è fra 0 e 1\,\%. È una buona notizia di natura
        pratica, perché significa che la prova definitiva del respiro può essere
        svolta a distanza realistica senza accorgimenti geometrici particolari.
  \item I canali \textbf{stazionari} danno 6,5--7,8 atti al minuto, troppo pochi per
        un adulto a riposo, e i gate prossimi al bersaglio sono saturi (gate~3 al
        100\,\%, gate~4 al 97\,\%).
\end{itemize}

\paragraph{Conclusione operativa.} Tre sessioni a distanze molto diverse --- 40~cm,
2,31~m e i circa 60~cm dello scenario sotto il banco (\cref{sec:sotto-banco}) ---
dicono la stessa cosa: \textbf{per il respiro vanno usati i canali in movimento}. Il
valore intorno a 7 atti al minuto dei canali stazionari è verosimilmente un artefatto
del filtraggio interno del canale e non respirazione --- ma, in assenza di ground
truth, questa resta un'ipotesi e va presentata come tale. A deciderlo è la prova con
metronomo descritta di seguito.

La terza sessione è anche l'unica svolta con un protocollo ripetuto: cinque trial
indipendenti nella geometria del progetto, con stima aggregata di
$21{,}0 \pm 2{,}5$ atti al minuto e un'analisi di sensibilità al criterio di
accettazione. È riportata in \cref{sec:sotto-banco} insieme al resto di quello
scenario, perché è lì che ha valore applicativo.

\subsection{La prova definitiva: respiro a ritmo imposto}
\label{sec:respiro-metronomo}

\emph{Acquisizione in corso.}

La misura valida per la tesi richiede una ground truth, e il modo più semplice di
ottenerla è imporre il ritmo: il soggetto respira seguendo un metronomo a
frequenza nota (per esempio 10, 15 e 20 atti al minuto), a distanza $\geq 1{,}5$~m,
per cinque trial per ritmo. Il confronto fra stima e ritmo imposto è la vera prova,
e permette di riportare un errore assoluto in atti al minuto invece di una
concordanza fra canali.

È inoltre necessario un \textbf{controllo negativo}: la stessa analisi applicata a
una registrazione di stanza vuota deve \emph{non} produrre un picco in banda
respiratoria. Senza questo controllo, un picco nello spettro non dimostra nulla, e
la banda 0,1--0,5~Hz è abbastanza larga da ospitare artefatti.

\section{Confronto fra i due moduli radar}
\label{sec:ld2420-risultati}

\emph{Attività parziale.}

Il LD2420 è stato collegato e verificato (\cref{sec:ld2420}): risponde ai movimenti,
riporta presenza e un valore di ``range'' in modalità ASCII a 115200 baud. In questa
modalità però non fornisce energia per-gate né distingue movimento e stazionario, e
l'unità del campo ``range'' non è documentata. Il confronto quantitativo con il
LD2410B richiede quindi la riconfigurazione del modulo in modalità binaria.

Alla data di questa stesura il confronto è limitato a un raffronto di specifiche
(\cref{tab:confronto-moduli}) e alla constatazione qualitativa che il LD2420
rileva presenza a distanze superiori. Il posizionamento dei due moduli resta quello
indicato dalle specifiche: il LD2410B è il modulo adatto al caso UPRISE (dettaglio
per-gate, respiro, distanze brevi), il LD2420 a una copertura di ambiente più ampia.

\section{Discussione}
\label{sec:discussione}

\subsection{Il quadro complessivo}

La \cref{tab:sintesi} raccoglie in una sola vista i risultati del capitolo. Tutte le
righe con soggetto in posizione H sono acquisite nella configurazione più favorevole
al PIR.

\begin{table*}[htbp]
\centering
\caption{Sintesi della campagna sperimentale. ``Rilev.'' è la frazione di campioni in
cui il sensore riporta la presenza, con la ground truth indicata.}
\label{tab:sintesi}
\small
\begin{tabularx}{\linewidth}{@{}llrrX@{}}
\toprule
\textbf{Scenario} & \textbf{Verità} & \textbf{Rilev.\ PIR} & \textbf{Rilev.\ radar} &
\textbf{Riferimento} \\
\midrule
Stanza vuota, 7,05~h & assente & $0$ ev/h & $0$ ev/h & \cref{sec:falsi-positivi} \\
Fermo seduto, 2,3~m & presente & $0{,}22\,\%$ & $\mathbf{100\,\%}$ & \cref{sec:persona-ferma} \\
Fermo in piedi, 1~m & presente & $1{,}52\,\%$ & $\mathbf{100\,\%}$ & \cref{sec:dose-risposta} \\
Micro-movimenti, 1~m & presente & $51{,}60\,\%$ & $\mathbf{100\,\%}$ & \cref{sec:dose-risposta} \\
Cammino sul posto, 1~m & presente & $85{,}20\,\%$ & $\mathbf{100\,\%}$ & \cref{sec:pir-movimento} \\
Cammino sul posto, 2--5~m & presente & $0{,}00\,\%$ & $\mathbf{100\,\%}$ & \cref{sec:pir-movimento} \\
Sotto il banco, immobile & presente & $1{,}32\,\%$ & $\mathbf{100\,\%}$ & \cref{sec:sotto-banco} \\
Sotto il banco, micro-mov. & presente & $96{,}52\,\%$ & $\mathbf{100\,\%}$ & \cref{sec:sotto-banco} \\
Due persone, ogni geometria & presenti & --- & $\mathbf{100\,\%}$ & \cref{sec:due-persone} \\
Vicino a 3~m, gate 2 & assente & $0{,}00\,\%$ & $0{,}00\,\%$ & \cref{sec:selettivita} \\
Vicino a 90\textdegree, gate 2 & assente & $0{,}00\,\%$ & $0{,}00\,\%$ & \cref{sec:selettivita} \\
\bottomrule
\end{tabularx}
\end{table*}

\subsection{Dove il mmWave vince}

\begin{enumerate}[leftmargin=1.8em,itemsep=3pt]
  \item \textbf{Sulla persona ferma il divario non è graduale, è categorico.} A 2,3~m
        seduto: $0{,}00\,\%$ contro $99{,}78\,\%$ di falsi negativi. A 1~m in piedi e
        nella configurazione migliore del PIR: $0{,}00\,\%$ contro $98{,}48\,\%$.
        Sotto il banco a 60~cm: $0{,}00\,\%$ contro $98{,}68\,\%$. Tre geometrie
        diverse, tre distanze diverse, lo stesso esito. Non si tratta di un sensore
        migliore di un altro, ma di un sensore che misura una grandezza a cui l'altro
        non ha accesso.
  \item \textbf{Il divario è attribuibile a una sola variabile.} La curva
        dose-risposta di \cref{sec:dose-risposta} varia la quantità di movimento a
        geometria, distanza e configurazione costanti, e produce $1{,}52\,\%$,
        $51{,}60\,\%$, $85{,}20\,\%$ sul PIR contro $100\,\%$ del radar in tutte e tre
        le condizioni. Il fallimento del PIR non dipende dalla distanza, dalla
        taratura o dal montaggio: dipende dall'immobilità del soggetto.
  \item \textbf{Il radar è più veloce anche dove il PIR dovrebbe essere avvantaggiato.}
        All'ingresso da una porta --- lo stimolo trasversale per cui la lente di
        Fresnel è progettata --- il radar rileva per primo in 10 trial su 10, con una
        differenza appaiata di $-0{,}60 \pm 0{,}35$~s.
  \item \textbf{Il radar aggiunge informazione, non solo affidabilità.} La distanza è
        misurata con linearità $R^2 = 0{,}99965$ e un errore di scala correggibile con
        un solo coefficiente; l'energia per-gate consente l'analisi del respiro; e
        l'energia del bersaglio e la dispersione della distanza forniscono due
        indicatori \emph{continui} e monotoni della quantità di movimento
        (\cref{tab:dose-risposta}), che sono la base dell'indice di vitalità
        (\cref{cap:vitalita}). Con un bit nessuna di queste grandezze è
        rappresentabile.
  \item \textbf{Sui falsi positivi i due sensori sono equivalenti.} Entrambi
        $\leq 0{,}43$ eventi/h su 7~h di osservazione, entrambi a zero eventi negli
        scenari di selettività. Il vantaggio del radar non si paga in affidabilità.
\end{enumerate}

\subsection{Dove il mmWave non vince}

Per equilibrio vanno riportati anche i punti a sfavore, che sono reali e che in un
caso sono peggiori di quanto le specifiche lascino prevedere.

\begin{itemize}[leftmargin=1.4em,itemsep=3pt]
  \item \textbf{Consumo}: tre ordini di grandezza in più (80~mA contro 50~\textmu A).
        È il vincolo che, da solo, determina l'architettura raccomandata nel
        \cref{cap:conclusioni}, e il \cref{cap:consumi} lo quantifica.
  \item \textbf{Il rilascio è più lento, e più lento del parametro dichiarato.} Il
        radar impiega $18{,}36 \pm 0{,}54$~s a dichiarare la stanza vuota contro i
        $12{,}96 \pm 1{,}40$~s del PIR: $5{,}40 \pm 0{,}67$~s in più, un effetto
        grande otto volte la propria incertezza. La coda propria del radar, stimata
        usando il PIR come cronometro, vale circa 8,9~s contro i \textbf{5~s} letti
        dal modulo come timeout di presenza. Il valore configurato non descrive il
        comportamento osservato, e tre osservazioni indipendenti concordano su 9--12~s.
  \item \textbf{Saturazione a distanza ravvicinata}: proprio nello scenario di
        interesse i canali più informativi saturano, e i canali stazionari per-gate
        non sono disponibili sotto 1,5~m. Sono vincoli gestibili --- sotto il metro si
        usano i canali in movimento, che risultano ricchi e correttamente collocati
        (\cref{sec:gate-sotto-banco}) --- ma vanno gestiti esplicitamente.
  \item \textbf{Sotto l'arredo la localizzazione è imprecisa}: la dispersione della
        distanza entro il trial è circa dieci volte peggiore che su un bersaglio
        immobile a 2,3~m. La presenza resta al $100\,\%$; è la distanza a non essere
        utilizzabile come misura fine in questa geometria.
  \item \textbf{Un bersaglio per canale}: il modulo non ha risoluzione angolare e
        riporta al più due bersagli, uno per canale. Separa bene due persone in stati
        diversi ($73{,}8\,\%$), male due persone ferme ($19{,}5\,\%$) e per nulla una
        persona dietro l'altra ($0{,}3\,\%$). Non è utilizzabile come contatore di
        presenze.
  \item \textbf{Complessità}: richiede UART a 256000 baud, la gestione di un
        protocollo a frame e --- per usarne i dati veri --- engineering mode e analisi
        numerica. Il PIR richiede un piedino digitale.
\end{itemize}

\subsection{La lettura per il progetto}

I dati portano a una conclusione che non è ``il PIR è un sensore scadente''. Il PIR,
configurato correttamente ed entro il proprio metro di portata utile, è un ottimo
rilevatore di \emph{movimento}: sotto il banco a 60~cm arriva al $96{,}5\,\%$, e lo fa
consumando tre ordini di grandezza in meno.

La conclusione corretta è che \textbf{il PIR fa bene un lavoro che non è questo}. La
grandezza che il progetto UPRISE deve misurare non è il movimento, è la
\emph{presenza di una persona che potrebbe non muoversi}: una persona intrappolata può
essere incosciente, ferita o esausta. In quella condizione il PIR è cieco per
principio fisico e non per taratura, come i dati mostrano in tre geometrie
indipendenti, e il mmWave la rileva nel $100\,\%$ dei campioni in ogni prova svolta.

Il corollario architetturale non è però ``sostituire il PIR''. I due sensori
falliscono e riescono in modo complementare, e il vincolo di consumo rende
impraticabile tenere il radar sempre acceso: l'architettura che i dati suggeriscono è
ibrida, con il PIR come sentinella a bassissimo consumo e il radar attivato in
modalità emergenza. È l'architettura discussa nel \cref{cap:conclusioni} e
quantificata nel \cref{cap:consumi}. La campagna sperimentale di questo capitolo ne
fornisce la giustificazione quantitativa: senza il radar, la persona immobile sotto il
banco non viene rilevata; senza il PIR, il nodo non ha autonomia.

\subsection{Limiti del setup sperimentale}
\label{sec:limiti}

I limiti da dichiarare esplicitamente sono i seguenti.

\begin{description}[leftmargin=1.4em,style=nextline,itemsep=3pt]
  \item[Numero di soggetti] Tutti i risultati comparativi sono stati ottenuti con un
    solo soggetto (indicato come S1 nel registro); negli scenari a due persone il
    secondo partecipante compare come bersaglio ma non è caratterizzato. La corporatura
    influisce sull'energia riflessa e, nel caso del PIR, sull'emissione infrarossa: la
    ripetizione degli scenari chiave con un secondo soggetto è la prima estensione da
    fare.
  \item[Ambiente] Stanza domestica, non un'aula scolastica: dimensioni, materiali,
    riflessioni multiple e presenza di arredi metallici sono diversi. Il trasferimento
    dei risultati al dimostratore reale va verificato.
  \item[Temperatura] Fra 25 e 29\,\celsius. Condizione sfavorevole al PIR in
    rilevamento: un confronto a 18--20\,\celsius\ è un'estensione a costo quasi nullo
    che rafforzerebbe il risultato, mostrando quanto della cecità sul fermo sia
    strutturale e quanto termico. La previsione, dalla fisica del principio, è che la
    cecità sul fermo resti totale a qualunque temperatura.
  \item[Range] 1--5~m per la caratterizzazione della distanza, per limite della stanza.
    Il caso UPRISE è sotto il metro e non ne è interessato.
  \item[Configurazione del PIR nella prima fase] Le acquisizioni fino al 22/08/2026
    sono state svolte con il ponticello di ritrigger in posizione L
    (\cref{sec:jumper}). Le serie sensibili a questa scelta sono state ripetute in H e
    sono quelle riportate; le altre sono state verificate insensibili per misura, non
    solo per deduzione. Resta il fatto che i due gruppi di serie sono separati da un
    rimontaggio del setup, che introduce uno scarto di alcuni centimetri sulla posizione
    del sensore: per questo le serie in H non sono state unite a quelle originali nella
    regressione della distanza.
  \item[Nessuna caratterizzazione angolare] La copertura angolare del \ldb\ non è stata
    misurata. Esiste un'osservazione qualitativa di rilevamento lateralmente arretrato e
    una misura puntuale a 90\textdegree\ con bersaglio fermo, ma non una curva. La
    misura è pianificata (soggetto a distanza fissa, angolo variato per passi noti) e
    non ancora svolta.
  \item[Vicino sempre fermo negli scenari di selettività] Il caso del vicino
    \emph{in movimento} a 90\textdegree\ non è stato provato ed è quello in cui il
    canale in movimento, libero, potrebbe agganciarlo.
  \item[Un solo esemplare per modello] Nessuna stima della variabilità fra esemplari.
    Il coefficiente di scala del $+3{,}81\,\%$ potrebbe essere specifico dell'esemplare
    in prova.
  \item[Assenza di strumentazione di misura] I consumi sono da datasheet
    (\cref{cap:consumi}) e le distanze sono riferite a marcature su pavimento, non a un
    distanziometro certificato. Questo impedisce, fra l'altro, di attribuire l'errore
    di scala del radar al modulo piuttosto che al riferimento.
  \item[Numero di ripetizioni disomogeneo] Gli scenari portanti hanno cinque o dieci
    trial, quelli di selettività e a due persone soltanto tre. Su tre trial la potenza
    statistica è bassa: le affermazioni corrispondenti sono formulate come assenza di
    effetto \emph{osservato}, non come assenza di effetto.
\end{description}

Nessuno di questi limiti mette in discussione il risultato centrale, che è robusto per
la natura stessa del divario misurato: la cecità del PIR sulla persona ferma è una
conseguenza del principio fisico, e i dati la confermano in tre geometrie indipendenti,
con cinque ripetizioni ciascuna, nella configurazione più favorevole al PIR.
